Of the eight redistricting reform successes recorded between 1990 and 2024 (P.0, Table~1), six were achieved through state ballot initiatives. The initiative channel bypasses the legislative collective action problem identified in P.0: voters can enact reform directly, without requiring incumbent legislators to vote against their interests. Where the legislature controls redistricting, the legislature controls its own reform; initiatives route around this circularity.

The initiative channel is available in 26 states---the 24 states without an initiative mechanism are entirely dependent on judicial or federal remedies (P.4, P.1). Among the 26 initiative states, redistricting reform has been attempted through the initiative process in at least 14, with a success rate of approximately 70\% for initiatives that actually reached the ballot.

Three states provide the most instructive case studies because their initiatives succeeded, were subsequently tested by legislative counter-measures, and generated substantial litigation and public documentation:

\begin{itemize}
  \item \textbf{California}: Propositions 11 (2008) and 20 (2010) created the Citizens Redistricting Commission for state legislative and congressional maps, respectively. Proposition 27 (2022) attempted to reverse these reforms and failed decisively (30.5\% yes). The California experience is the longest-running and most tested example of initiative-based redistricting reform.

  \item \textbf{Colorado}: Amendments Y (congressional) and Z (state legislative) passed in 2018 with 71\% and 70\% support, creating independent commissions for both chambers. Colorado's simultaneous treatment of congressional and legislative redistricting---and the commissions' extensive 2021 redistricting process---provides a detailed recent case study.

  \item \textbf{Michigan}: Proposal 2 (2018) created the Michigan Independent Citizens Redistricting Commission (MICRC) and passed with 61\% support. Michigan's 2021 redistricting was notable for extensive public engagement and for producing maps that were subsequently challenged but upheld in federal court.
\end{itemize}

This paper's contribution is not merely descriptive. We extract the design principles that made these initiatives durable and successful, and we apply them to develop model ballot initiative language that would add an algorithmic redistricting requirement to an existing commission structure. Section~2 analyzes each case study in turn. Section~3 distills common design elements. Section~4 presents model language for California. Section~5 concludes.
